\chapter{Task 3 - Temperature and Humidity Monitoring with OLED Display}

This chapter details the design and implementation of the environmental monitoring subsystem, which integrates an OLED display for real-time data visualization. The primary objective of this task is to provide immediate, intuitive feedback on ambient conditions to the user. We designed the system to continuously monitor temperature and humidity levels and dynamically adapt the visual interface based on predefined operational states: Normal, Warning, Critical, and Fire Alert. This adaptive display strategy ensures that critical safety information is prioritized over routine data, thereby enhancing the system's usability and responsiveness in emergency scenarios.

\section{OLED Layout Modes}

We utilized an SSD1306 OLED display (128x64 pixels) to present environmental data. Rather than employing a static layout, we implemented a dynamic rendering engine that switches between four distinct display modes based on the current system state. This approach allows the system to highlight the most relevant information for each specific condition.

\subsection{Normal State}
In the \textbf{Normal} state (Temperature 15--30\textdegree C, Humidity 40--80\%), the display is configured to maximize readability of the sensor readings. We designed a clean, uncluttered layout that prominently features the current temperature and humidity values. To provide immediate reassurance to the user, the interface includes a comfort status indicator (displaying "IDEAL" or "GOOD") and a graphical smiley icon. This mode is intended for routine monitoring when environmental conditions are within safe limits.

\begin{figure}[H]
    \centering
    \includegraphics[width=0.6\textwidth]{assets/oled_normal.png}
    \caption{OLED display in Normal state showing temperature, humidity, and comfort icon.}
    \label{fig:oled_normal}
\end{figure}

\subsection{Warning State}
The \textbf{Warning} state is triggered when environmental parameters drift outside the optimal comfort range (Temperature \textless 15\textdegree C or \textgreater 33\textdegree C, or Humidity \textless 40\% or \textgreater 80\%). To attract the user's attention, we implemented a visual blinking effect on the display border with a 500ms interval. The layout shifts focus from detailed readings to specific warning messages (e.g., "HOT", "COLD", "DRY", "TOO HUMID"), accompanied by an exclamation icon. This visual hierarchy ensures that the user is promptly alerted to the deviation before conditions deteriorate further.

\subsection{Critical State}
In the \textbf{Critical} state (Temperature \textgreater 40\textdegree C or Humidity \textless 30\%), the system escalates the visual urgency. We inverted the display header colors and implemented a flashing effect to maximize contrast and visibility. The screen presents the sensor values alongside imperative corrective action messages (e.g., "TOO HOT! Cool down!" or "TOO DRY! Add moisture!"). This design is intended to guide the user toward immediate remediation steps to prevent equipment damage or health risks.

\subsection{Fire Alert State}
The \textbf{Fire Alert} state represents the highest priority condition in the system hierarchy. Upon detection of a flame by the sensor, the display logic overrides all other data to present a full-screen emergency alert. We designed this mode to display "FIRE DETECTED!" and "EVACUATE NOW" in large, bold text, ensuring that the warning is legible from a distance. In this state, routine sensor data is suppressed to prevent distraction during an emergency.

\begin{figure}[H]
    \centering
    \includegraphics[width=0.6\textwidth]{assets/oled_fire.png}
    \caption{OLED display in Fire Alert state showing emergency evacuation message.}
    \label{fig:oled_fire}
\end{figure}

\section{State Selection Flow}

To manage the transitions between these display modes, we implemented a semaphore-based control mechanism within the FreeRTOS environment. Instead of relying on complex conditional logic that polls global variables, the display task waits for specific binary semaphores that signal the current system state. This design ensures that the visual output remains strictly synchronized with the core system logic.

The state evaluation logic enforces a strict priority hierarchy: \textbf{Fire Alert} takes precedence over \textbf{Critical}, which overrides \textbf{Warning}, with \textbf{Normal} being the default state. This hierarchical approach guarantees that the most urgent information is always displayed, regardless of other concurrent conditions.

The following pseudocode illustrates the selection flow implemented in our display task:

\begin{verbatim}
FUNCTION updateDisplay():
    Acquire I2C Mutex
    
    // Check which state is active
    state = getSystemState()
    
    SWITCH state:
        CASE FIRE_ALERT:
            drawFireAlertDisplay(animationFrame)
            BREAK
            
        CASE CRITICAL:
            drawCriticalDisplay(temp, humi, animationFrame)
            BREAK
            
        CASE WARNING:
            drawWarningDisplay(temp, humi, animationState)
            BREAK
            
        CASE NORMAL:
            drawNormalDisplay(temp, humi)
            BREAK
            
    Release I2C Mutex
    
    // Adjust refresh rate based on urgency
    IF state >= WARNING:
        Delay 300ms  // Fast refresh for alerts
    ELSE:
        Delay 1000ms // Save power in normal mode
END FUNCTION
\end{verbatim}

\section{Queue and Semaphore Data Flow}

A key architectural requirement for this project was the elimination of global variables to ensure thread safety and modularity. We achieved this by encapsulating all sensor readings within a \texttt{SensorData\_t} structure, which is transmitted between tasks via FreeRTOS queues. Furthermore, system states are managed exclusively through binary semaphores.

Our implementation removes global variable dependencies through three specific mechanisms:

\begin{enumerate}
    \item \textbf{Queue-Based Data Transfer:} All sensor values are packed into the \texttt{SensorData\_t} struct and passed through a queue. This prevents race conditions associated with shared global variables.
    \item \textbf{Semaphore-Based State Signaling:} We utilize binary semaphores to signal state changes, providing a robust synchronization primitive that replaces global state flags.
    \item \textbf{Inter-Task Communication:} Tasks communicate solely through these OS primitives, ensuring that data access is thread-safe and deterministic.
\end{enumerate}

The OLED task structure demonstrates this pattern. It initializes the display and sensor, then enters an infinite loop where it reads data, evaluates the system state, updates semaphores, and renders the appropriate display layout. The state evaluation is implemented as a pure function, ensuring no side effects, while the state update function uses mutexes to guarantee atomic transitions between states.

\begin{verbatim}
TASK temp_humi_oled():
    Initialize OLED display (with I2C mutex protection)
    Initialize DHT11 sensor
    
    WHILE TRUE:
        Read temperature and humidity from DHT11
        Validate readings
        
        // Pack data into struct (no global variables)
        sensorData.temperature = temperature
        sensorData.humidity = humidity
        
        // Evaluate state using pure function
        newState = evaluateSystemState(temp, humidity, flame)
        
        // Update semaphores (thread-safe)
        updateSystemState(newState)
        
        // Send to queue for other tasks
        sendSensorData(sensorData)
        
        // Render display based on state
        Acquire I2C mutex
        currentState = getSystemState()
        SWITCH currentState:
            CASE FIRE_ALERT: drawFireAlertDisplay()
            CASE CRITICAL: drawCriticalDisplay()
            CASE WARNING: drawWarningDisplay()
            CASE NORMAL: drawNormalDisplay()
        Release I2C mutex
        
        // Adaptive delay based on urgency
        IF state >= WARNING: Delay 300ms
        ELSE: Delay 1000ms
    END WHILE
END TASK
\end{verbatim}

This architecture ensures that the display task operates independently yet synchronously with the rest of the system, providing reliable visual feedback without compromising the stability of the real-time operating system.
